\textbf{9.} Докажите, что задача поиска изоморфизма между графами сводится по Куку к задаче распознавания, существует ли изоморфизм.

\textbf{Решение.} Первое, что приходит в голову, --- давайте строить изоморфизм постепенно, беря из каждого графа по вершине. А именно, зафиксируем какую-то вершину $i$ в первом графе $G_1$ и будем циклом перебирать вершину $j$ во втором графе $G_2$. Хотим узнать, есть ли изоморфизм, переводящий $i$ в $j$. Если их степени не совпадают, то перебираем дальше -- нужного изоморфизма нет. Иначе прицепим к этим вершинам <<метлу>> --- $n - deg(i)$ ($ = n - deg(j)$) вспомогательных вершин (см. рисунок).
\begin{center}
    \begin{asy}
size(220,0);

unitsize(2cm);
texpreamble("\usepackage[T2A]{fontenc}\usepackage[utf8]{inputenc}
\usepackage{mathtext}\usepackage[russian]{babel}");

draw((1, 1)--(-1, 1)--(-1, 0)--(1, 0)--cycle);
draw((4, 1)--(2, 1)--(2, 0)--(4, 0)--cycle);
label("$G_1$", (0, 0.5), N);
label("$G_2$", (3, 0.5), N);

draw(shift(0, 0.25)*scale(1/4)*unitcircle);
draw(shift(3, 0.25)*scale(1/4)*unitcircle);
label("$i$", (0, 0.25));
label("$j$", (3, 0.25));

draw((0, 0)--(-0.75, -0.25));
path p = shift(-0.75, -0.25)*scale(0.125)*unitcircle;
draw(p);
filldraw(p, white);
draw((0, 0)--(-0.5, -0.5));
p = shift(-0.5, -0.5)*scale(0.125)*unitcircle;
draw(p);
filldraw(p, white);
draw((0, 0)--(0.75, -0.25));
p = shift(0.75, -0.25)*scale(0.125)*unitcircle;
draw(p);
filldraw(p, white);
draw((0, 0)--(0.5, -0.5));
p = shift(0.5, -0.5)*scale(0.125)*unitcircle;
draw(p);
filldraw(p, white);
label("$\ldots$", (0, -0.75));
draw((0, 0)--(-0.2, -0.6));
draw((0, 0)--(0.2, -0.6));
        
draw((3, 0)--(2.25, -0.25));
path p = shift(2.25, -0.25)*scale(0.125)*unitcircle;
draw(p);
filldraw(p, white);
draw((3, 0)--(2.5, -0.5));
p = shift(2.5, -0.5)*scale(0.125)*unitcircle;
draw(p);
filldraw(p, white);
draw((3, 0)--(3.75, -0.25));
p = shift(3.75, -0.25)*scale(0.125)*unitcircle;
draw(p);
filldraw(p, white);
draw((3, 0)--(3.5, -0.5));
p = shift(3.5, -0.5)*scale(0.125)*unitcircle;
draw(p);
filldraw(p, white);
label("$\ldots$", (3, -0.75));
draw((3, 0)--(2.8, -0.6));
draw((3, 0)--(3.2, -0.6));

draw((1.5, -1)--(1, -0.6), EndArrow);
draw((1.5, -1)--(2, -0.6), EndArrow);
label("по $n - deg(i)$ вершин", (1.5, -1), S);
    \end{asy}
\end{center}

Проверим с помощью оракула $\mathsf{GI}$, изоморфны ли полученные графы. Если да, то вершине $i$ может соответствовать только $j$, потому что только у них двоих степень вершины равна $n$. Но тогда элементы метлы, как вершины со степенью 1, перешли в такие же листы. Поэтому если убрать $n$ листов у $i$ и соответствующие им листы у $j$, получится изоморфизм $G_1$ и $G_2$, который переводит $i$ в $j$. Обратно, если такой изоморфизм есть, то его можно легко продолжить до изоморфизма построенных графов переводом листов метлы $i$ в листы метлы $j$.

Как только найдём $j$, для которого построенные графы изоморфны, мы запоминаем пару $(i, j)$ и начинаем перебирать новые: фиксируем другую вершину $i' \in V(G_1) \setminus \{i\}$ и в цикле проходимся по $j' \in V(G_2) \setminus \{j\}$. Теперь хотим проверить, есть ли изоморфизм, переводящий $i$, $i'$ в $j$, $j'$ соответственно. После проверки на одинаковость степеней совершаем действия, аналогичные выше: к $i$ и $j$ прицепляем метлу размера $n - deg(i)$, а к $i'$ и $j'$ -- размера $n - deg(i') + 1$. По тем же соображениям между исходными графами есть нужный изоморфизм тогда и только тогда, когда он есть между новыми, причём мётлы будут переходить в мётлы, а из-за того, что степень $i$, $j$ и $i'$, $j'$ теперь равна $n$ и $n+1$ соответственно, изоморфизм, если он есть, обязан перевести $i$ в $j$, а $i'$ в $j'$. 

В общем случае, если уже найдены вершины $i_1, \ldots, i_r$ и $j_1, \ldots, j_r$ такие, что существует изоморфизм $\phi\colon V(G_1) \to V(G_2)$, $\phi(i_k) = j_k$, то циклом перебираются кандидаты на роль $i_{r+1}$ и $j_{r+1}$. При рассмотрении очередной такой пары к вершинам $i_k$ и $j_k$ прицепляется метла размером $n - deg(i_k) + k$ и проверяется наличие изоморфизма между полученными графами. По соображениям выше это рано или поздно даст нам нужную пару. Каждая итерация делается за полином от размера графа.

P. S. Решение этой задачи стоило Илье Данииловичу 10 хинкалей...